We start by motivating our method before explaining its details in \Cref{subsec:desc_algo}. Many successful self-supervised learning approaches build upon the cross-view prediction framework introduced in~\cite{becker1992self}. Typically, these approaches learn representations  by predicting different views (e.g., different random crops) 
of the same image from one another. Many such  approaches cast the prediction problem directly in representation space: the representation of an  \augview~of an image should be predictive of the representation of another \augview~of the same image. However, predicting directly in representation space can lead to collapsed representations: for instance, a representation that is constant across views is always fully predictive of itself. Contrastive methods circumvent this problem by reformulating the prediction problem into one of discrimination: from the representation of an \augview, they learn to discriminate between the representation of another \augview~of the same image, and the representations of \augview s of different images. In the vast majority of cases, this prevents the training from finding collapsed representations. Yet, this discriminative approach typically requires comparing each representation of an \augview~with many negative examples, to find ones sufficiently close to make the discrimination task challenging. In this work, we thus tasked ourselves to find out whether these negative examples are indispensable to prevent collapsing while preserving high performance.

To prevent collapse, a straightforward solution is to use a fixed randomly initialized network to produce the targets for our predictions. While avoiding collapse, it empirically does not result in very good representations. Nonetheless, it is interesting to note that the representation obtained using this procedure can already be much better than the initial fixed representation. In our ablation study (\Cref{sec:ablation}), we apply this procedure by predicting a fixed randomly initialized network and achieve $18.8\%$ top-1 accuracy (\Cref{tab:ablation_bootstrap}) on the linear evaluation protocol on ImageNet, whereas the randomly initialized network only achieves $1.4\%$ by itself. This experimental finding is the core motivation for \algo: from a given representation, referred to as \textit{target}, we can train a new, potentially enhanced representation, referred to as \textit{online}, by predicting the target representation. From there, we can expect to build a sequence of representations of increasing quality by iterating this procedure, using subsequent online networks as new target networks for further training. In practice, \algo generalizes this bootstrapping procedure by iteratively refining its representation, but using a slowly moving exponential average of the online network as the target network instead of fixed checkpoints.

\subsection{Description of \algo}
\label{subsec:desc_algo}
\begin{figure}[!ht]
\vskip -35em
    \centering
    %\includestandalone[mode=buildnew]{algo}
    %\includegraphics[scale=0.75, bb=0 0  461 1233]{img/fadingdirtyfix.png}
        \includegraphics[bb=0 0  1233 461]{img/archi.pdf}
    \caption{\algo's architecture. \algo minimizes a similarity loss between $q_\netparams(z_\netparams)$ and $\mathrm{sg}(z'_\targetparams)$, where $\netparams$ are the trained weights, $\targetparams$ are an exponential moving average of $\netparams$ and $\mathrm{sg}$ means stop-gradient. At the end of training, everything but $f_\netparams$ is discarded, and $y_\netparams$ is used as the image representation. }
    \label{fig:algo}
    %\vskip -1em
\end{figure}
\algo's goal is to learn a representation $y_\netparams$ which can then be used for downstream tasks. As described previously, \algo uses two neural networks to learn: the \emph{online} and \emph{target} networks. The online network is defined by a set of weights $\netparams$ and is comprised of three stages: an \emph{encoder} $f_\netparams$, a \emph{projector}~$g_\netparams$ and a \emph{predictor}~$q_\netparams$, as shown in \Cref{fig:algo} and \Cref{fig:BYOL_sketch}. The target network has the same architecture as the online network, but uses a different set of weights $\targetparams$.
The target network provides the regression targets to train the online network, and its parameters $\targetparams$ are an exponential moving average of the online parameters $\netparams$~\cite{DDPG}. More precisely, given a target decay rate $\tau \in [0, 1]$, after each training step we perform the following update,
\begin{equation}
\targetparams \leftarrow \tau\targetparams + (1 - \tau) \netparams.  \label{eq:tau}  
\end{equation}
Given a set of images $\mathcal{D}$, an image $x \sim \mathcal{D}$ sampled uniformly from $\mathcal{D}$, and two distributions of \imageaugmentation s $\mathcal{T}$ and $\mathcal{T}'$, \algo produces two \augview s $v \eqdef t(x)$ and $v'\eqdef t'(x)$ from $x$ by applying respectively \imageaugmentation s $t\sim \mathcal{T}$ and $t'\sim \mathcal{T}'$. From the first \augview~$v$, the online network outputs a \emph{representation} $y_\netparams \eqdef f_\netparams(v)$ and a projection $z_\netparams \eqdef g_\netparams(y)$. The target network outputs $y'_\targetparams \eqdef f_\targetparams(v')$ and the \emph{target projection} $z'_\targetparams \eqdef g_\targetparams (y')$ from the second \augview~$v'$. We then output a \emph{prediction} $q_\netparams(z_\netparams)$ of $z'_\targetparams$ and $\ell_2$-normalize both $q_\netparams(z_\netparams)$ and $z'_\targetparams$ to $\bar{q_\netparams}(z_\netparams) \eqdef q_\netparams(z_\netparams)/\normtwo{q_\netparams(z_\netparams)}$ and $\bar{z}'_\targetparams \eqdef z'_\targetparams / \|z'_\targetparams\|_2$.  Note that this predictor is only applied to the online branch, making the architecture asymmetric between the online and target pipeline. Finally we define the following mean squared error between the normalized predictions and target projections,\footnote{While we could directly predict the representation~$y$ and not a projection $z$, previous work~\cite{SimCLR} have empirically shown that using this projection improves performance.}
\begin{equation}
    \mathcal{L}_{\netparams, \targetparams} \eqdef \normtwo{\bar{q_{\theta}}(z_\netparams) - \bar{z}'_\targetparams}^2 = 2 - 2\cdot\frac{\langle q_\netparams(z_\netparams),  z'_\targetparams \rangle }{\big\|q_\netparams(z_\netparams)\big\|_2\cdot
    \big\|z'_\targetparams\big\|_2
    }
    \cdot
    \label{eq:cosine-loss}
\end{equation}

We symmetrize the loss $\mathcal{L}_{\netparams, \targetparams}$ in \Cref{eq:cosine-loss} by separately feeding $v'$ to the online network and $v$ to the target network to compute $\tilde{\mathcal{L}}_{\netparams, \targetparams}$. At each training step, we perform a stochastic optimization step to minimize $\mathcal{L}^\algo_{\netparams, \targetparams} = \mathcal{L}_{\netparams, \targetparams} + \tilde{\mathcal{L}}_{\netparams, \targetparams}$ with respect to $\netparams$ only, but \textit{not} $\targetparams$, as depicted by the stop-gradient in \Cref{fig:algo}. \algo's dynamics are summarized as
\begin{align}
&\netparams \leftarrow \mathrm{optimizer} \left(\netparams, \nabla_\netparams \mathcal{L}^\algo_{\netparams, \targetparams}, \eta \right),
\\
&\targetparams \leftarrow \tau \targetparams + (1 - \tau) \netparams,  \tag{\ref{eq:tau}}
\end{align}
where $\mathrm{optimizer}$ is an optimizer and $\eta$ is a learning rate.

At the end of training, we only keep the encoder $f_\netparams$; as in~\cite{MOCO}. When comparing to other methods, we consider the number of inference-time weights only in the final representation $f_\netparams$.
The full training procedure is summarized in \Cref{app:algo}, and \texttt{python} pseudo-code based on the libraries \texttt{JAX}~\cite{jax2018github} and \texttt{Haiku}~\cite{haiku2020github} is provided in
in Appendix~\hyperlink{page.33}{J}.
%; \Cref{app:pseudocode}. 

\subsection{Intuitions on \algo's behavior}
\label{sec:intuition_byol}
As \algo does not use an explicit term to prevent collapse (such as negative examples~\cite{CPC}) while minimizing $\mathcal L_{\netparams,\targetparams}^\algo$ with respect to $\netparams$, it may seem that \algo should converge to a minimum of this loss with respect to $(\netparams, \targetparams)$ (\emph{e.g.}, a collapsed constant representation). However \algo's target parameters $\targetparams$ updates are \textbf{not} in the direction of $\nabla_\targetparams \mathcal L_{\netparams,\targetparams}^\algo$. More generally, we hypothesize that there is no loss $L_{\netparams, \targetparams}$ such that \algo's dynamics is a gradient descent on $L$ jointly over $\netparams, \targetparams$. This is similar to GANs~\cite{goodfellow2014generative}, where there is no loss that is jointly minimized w.r.t.\,both the discriminator and generator parameters. 
There is therefore no \textit{a priori} reason why \algo's parameters would converge to a minimum of $\mathcal{L}^\algo_{\netparams, \targetparams}$. 

While \algo's dynamics still admit undesirable equilibria, we did not observe convergence to such equilibria in our experiments. In addition, when assuming \algo's predictor to be optimal\footnote{For simplicity we also consider \algo without normalization (which performs reasonably close to \algo, see \Cref{app:abla_norm}) nor symmetrization} i.e., $q_\netparams = q^\star$ with 

\begin{equation}
\label{eq:optimal_predictor}
    q^\star \eqdef \argmin_q \EE{\normtwo{q(z_\netparams) - z'_\targetparams}^2}, \quad \text{where} \quad q^\star(z_\netparams) = \EE{z'_\targetparams | z_\netparams},
\end{equation}

we hypothesize that the undesirable equilibria are unstable. Indeed, in this optimal predictor case, \algo's updates on~$\netparams$ follow in expectation the gradient of the expected conditional variance (see \Cref{app:thm-envelope} for details),
\begin{equation}
\label{eq:byol_minimize_cond_var}
    \nabla_\netparams\EE{\normtwo{q^\star(z_\netparams) - z'_\targetparams}^2} = \nabla_\netparams\EE{\normtwo{\EE{z'_\targetparams | z_\netparams} -z'_\targetparams }^2} = \nabla_\netparams \EE{ \sum_i \var(z'_{\targetparams, i} | z_\netparams)},
\end{equation}
where $z'_{\targetparams, i}$ is the $i$-th feature of $z'_{\targetparams}$. 

Note that for any random variables $X,$ $Y,$ and $Z$, $\var(X|Y,Z) \leq \var(X|Y)$. Let $X$ be the target projection, $Y$ the current online projection, and $Z$ an additional variability on top of the online projection induced by stochasticities in the training dynamics: purely discarding information from the online projection cannot decrease the conditional variance. 

In particular, \algo avoids constant features in $z_\netparams$ 
as, for any constant $c$ and random variables $z_\netparams$ and $z'_\targetparams$, $\var(z'_\targetparams | z_\netparams) \leq \var(z'_\targetparams | c)$;  hence our hypothesis on these collapsed constant equilibria being unstable. Interestingly, if we were to minimize $\mathbb{E}[\sum_i \var(z'_{\targetparams, i} | z_\netparams)]$ with respect to $\targetparams$, we would get a collapsed $z'_\targetparams$ as the variance is minimized for a constant $z'_\targetparams$. 
Instead, \algo makes $\targetparams$ closer to $\netparams$, incorporating sources of variability captured by the online projection into the target projection. 

Furthemore, notice that performing a hard-copy of the online parameters $\netparams$ into the target parameters $\targetparams$ would be enough to propagate new sources of variability. However, sudden changes in the target network might break the assumption of an optimal predictor, in which case \algo's loss is not guaranteed to be close to the conditional variance. We hypothesize that the main role of \algo's moving-averaged target network is to ensure the near-optimality of the predictor over training; \Cref{sec:ablation} and \Cref{app:importance_no_pred} provide some empirical support of this interpretation.

\subsection{Implementation details}
\label{sec:impl_detail}
\paragraph{\Imageaugmentation s} \algo uses the same set of \imageaugmentation s as in \simclr~\cite{SimCLR}. First, a random patch of the image is selected and resized to $224 \times 224$ with a random horizontal flip, followed by a color distortion, consisting of a random sequence of brightness, contrast, saturation, hue adjustments, and an optional grayscale conversion. Finally Gaussian blur and solarization are applied to the patches. Additional details on the \imageaugmentation s are in \Cref{app:transf}.
\paragraph{Architecture} We use a convolutional residual network \cite{ResNet} with 50 layers and post-activation (ResNet-$50(1 \times)$ v1) as our base parametric encoders $f_\netparams$ and $f_\targetparams$. We also use deeper ($50$, $101$, $152$ and $200$ layers) and wider (from $1 \times$ to $4 \times$) ResNets, as in \cite{zagoruyko2016wide,kolesnikov2019revisiting,SimCLR}. Specifically, the representation $y$ corresponds to the output of the final average pooling layer, which has a feature dimension of $2048$ (for a width multiplier of $1\times$). As in \simclr~\cite{SimCLR}, the representation~$y$ is projected to a smaller space by a \emph{multi-layer perceptron} (MLP)~$g_\netparams$, and similarly for the target projection $g_\targetparams$. This MLP consists in a linear layer with output size $4096$ followed by batch normalization~\cite{BatchNorm}, rectified linear units (ReLU)~\cite{nair2010rectified}, and a final linear layer with output dimension $256$. Contrary to \simclr, the output of this MLP is not batch normalized. The predictor $q_\netparams$ uses the same architecture as $g_\netparams$.

\paragraph{Optimization}
We use the \LARS optimizer~\cite{LARS} with a cosine decay learning rate schedule~\cite{SGDR}, without restarts, over $1000$ epochs, with a warm-up period of $10$ epochs. We set the base learning rate to $0.2,$ scaled linearly~\cite{Goyal2017} with the batch size ($\text{LearningRate} = 0.2\times \text{BatchSize} / 256$). In addition, we use a global weight decay parameter of $1.5\cdot10^{-6}$ while excluding the biases and batch normalization parameters from both \LARS adaptation and weight decay. For the target network, the exponential moving average parameter $\tau$ starts from $\tau_\text{base} = 0.996$ and is increased to one during training. Specifically, we set $\tau \triangleq 1 - (1-\tau_\text{base})\cdot \pa{\cos\pa{\pi k / K} + 1} / 2$ with $k$ the current training step and $K$ the maximum number of training steps. We use a batch size of $4096$ split over $512$ Cloud TPU v$3$ cores. With this setup, training takes approximately $8$ hours for a ResNet-$50(\times 1)$. All hyperparameters are summarized in Appendix~\hyperlink{page.33}{J}; 
%\Cref{code:hyperparams}; 
an additional set of hyperparameters for a smaller batch size of $512$ is provided in \Cref{app:smaller-batch}.